\chapter{Provenance and claim boundary}\label{chap:provenance}

\section{The source and the object of study}

The object of this transcript is a single section of the paper of Buchholz and
MacDougall in J.\ Number Theory \textbf{81} (2000)~\cite{BuchholzMacDougall2000};
specifically \S2.2.4, ``Vertical translation of a quartic'' (l.527--632 of the
verbatim layout, md5 \texttt{a177e2e673705c0d68cc5fdf3a6626ce}).  That section supplies the last
inferential step of the paper's Theorem~7, the classification of rational-derived
quartics.  A quartic is \emph{rational-derived} when it and all of its derivatives
split into linear factors over $\Q$; Theorem~7 asserts that every rational-derived
quartic is equivalent, modulo the affine group $\langle X^{*}\rangle$, to one of an
explicitly parametrised family, the last clause of which is governed by the elliptic
curve $E\colon z^{2}=w(w-6)(w+18)$ (l.630).

The transcript is a survey of one section, not a verdict on the paper.  Buchholz and
MacDougall's verified results (their Theorems~5 and~6, the identity
$\Delta(f')=-16\,D_{6}(a,b)$ at l.606, and the discriminant tower at l.568) remain
machine-verified true; typographical errors and inferential gaps are distinct
phenomena and are kept distinct throughout.

\section{The promise (T0) and the assumption (T2), quoted}

The section opens with its promise.  It sets up a rational-derived quartic, moves the
highest local minimum up to become a double root by a rational vertical translate, and
undertakes to show that the two remaining roots are rational (l.529--534):

\begin{quote}
``If we now allow our rational-derived quartics to undergo a vertical translate by a
rational distance, so that the (highest) local minimum, $x_{\min}$ say, is moved up to
become a double root, then the remaining two roots, $r$ and $s$ say, of the quartic
could possibly lie in a quadratic extension of $\Q$. Certainly, we have no a priori
reason to expect the extra roots to be rational (see Figure~2). None-the-less, in this
section we show that the latter is precisely the case.'' \hfill (l.529--534)
\end{quote}

This is the terminus of the transcript: the proposition the section undertakes to
prove.  In the formalisation it is the object \lean{BMThm7Transcript.T0_claim}
(Chapter~\ref{chap:transcript}).

The argument then names its standing datum.  Having introduced
$f(x)=x(x-1)(x-a)(x-b)$ (l.549--551) and a rational translate amount $c$, it writes
(l.553--558):

\begin{quote}
``We assume that the resulting quartic, $F(x)$ say, given by
\[
F(x) = x(x-1)(x-a)(x-b) + c \tag{4}
\]
is rational-derived and has a double root. As before, we require $\Delta(F)=0$\dots''
\hfill (l.553--558)
\end{quote}

The verb is \emph{assume}.  The clause ``$F$ is rational-derived'' is, under the
standing hypothesis that $f$ is rational-derived, exactly the section's own
conclusion; this is the pre-registered location of the gap
(P1, Chapter~\ref{chap:transcript}).  In the formalisation the assumption is
\lean{BMThm7Transcript.T2_assumption}.

\section{Claim boundary}

The transcript establishes nothing for or against Theorem~7 over $\Q$, and by
construction it cannot.  Theorem~7 over $\Q$ is \emph{not} refuted here: a
counterexample over $\Q$ would refute Conjecture~1 of the same paper, which is open
(RESULTS \S7).  What the transcript does establish is a boundary result: the printed
argument of \S2.2.4 does not derive its conclusion from the results preceding it, and
every known route to the conclusion passes through Conjecture~1
(Chapter~\ref{chap:boundary}).  The two refutations carried out in the formalisation
are refutations of the printed \emph{implication} over $\Q$
(\lean{BMThm7Transcript.not_T9_bridge_Q}) and of the \emph{statement} over a specific
quadratic field $K=\Q(\sqrt{3441})$ (\lean{BMThm7Transcript.not_T0_claim_K}); neither
is a statement about Theorem~7 over $\Q$.

\section{Pre-registration discipline}

The segmentation of \S2.2.4 into inferential units $\mathrm{T0}$--$\mathrm{T10}$, the
classification of each unit, and seven numbered predictions $\mathrm{P1}$--$\mathrm{P7}$
about where the argument would and would not formalise, were fixed in a document
(STEPS.md) whose md5 was frozen at \texttt{3ccf4bb1620b264baece2e2e510ad283} at
12:25 on 2026-07-16c, \emph{before} any Lean build was attempted.  The freeze precedes
every verdict; the adjudication record (ADJUDICATION.md) postdates it in full.  The
discipline is that predictions are compared against the kernel and recorded as met or
unmet, and are not rewritten after the fact.  All seven predictions were confirmed;
two (P5 and~P7) were confirmed in a form stronger than predicted, and one incidental
finding upgraded the symmetric-case lemma T7a to an in-house proof
(Chapter~\ref{chap:transcript}).

The blueprint format below follows the leanblueprint convention used by projects such
as PFR~\cite{pfr-blueprint}, the Equational Theories Project~\cite{equational-theories-blueprint},
and PrimeNumberTheoremAnd~\cite{pnt-blueprint}, from which the template plumbing of this
document was adapted.
